\section{Background}

This section introduces the foundations of the MouseBrainBench framework. First,
public mouse-brain resources are addressed as the empirical basis of the
proposal. Second, visual prediction benchmarks are reviewed because they
represent one of the most active bridges between neuroscience and Artificial
Intelligence. Third, connectomic resources are considered, paying special
attention to the difference between local high-resolution evidence and
whole-brain claims. Finally, existing simulation tools and the notion of
mechanistic identifiability are discussed.

\subsection{Public mouse-brain resources}

The Allen Brain Observatory and Allen Visual Behavior Neuropixels have promoted
the reuse of large-scale recordings from the mouse visual system
\citep{deVries2022AllenSharing,siegle2021survey}. These resources are relevant
because they provide standardized data that can be processed by different
research groups and algorithms. Therefore, they are especially useful for
evaluating reproducibility, population responses, and predictive models.

However, these recordings do not automatically solve the mechanistic
identifiability problem. A model can reproduce a response pattern observed in
these data without proving that its internal structure corresponds to the
biological circuit that produced the activity. For this reason,
MouseBrainBench uses Allen data as a negative and conservative control. If the
framework is well designed, it must be able to detect reproducible information
while rejecting unsupported anatomical or causal interpretations.

\subsection{Visual prediction benchmarks}

Predictive models of the mouse visual system have gained attention because they
allow computational approaches from Artificial Intelligence to be tested against
neural responses. Sensorium introduced a standardized benchmark for static
visual input using large-scale recordings from mouse primary visual cortex
\citep{willeke2022sensorium}. Dynamic Sensorium extended this idea to video
stimuli and out-of-distribution evaluation \citep{turishcheva2023dynamic,
turishcheva2024retrospective}.

These benchmarks are important because they provide common data, baselines, and
evaluation protocols. Nevertheless, predictive quality should not be confused
with mechanistic explanation. A model that predicts well may be useful as a
functional surrogate, but it does not necessarily identify the anatomical
structure or causal process behind the response. MouseBrainBench therefore
treats Sensorium as a predictive layer and not as direct evidence of a digital
twin.

\subsection{Connectomics and structure-function resources}

Connectomics provides anatomical information about neural circuits, but its
interpretation depends strongly on scale. Local electron-microscopy resources
can provide synaptic-level information in restricted volumes, whereas whole-brain
resources usually operate at coarser spatial resolution. MICRONS is a key
resource because it links functional recordings and electron-microscopy
reconstruction in mouse visual cortex \citep{bae2025microns}. This makes it
suitable for evaluating local structure-function relationships.

Even in this case, the resource must be interpreted carefully. MICRONS is large
and biologically rich, but it is not a complete synaptic reconstruction of the
whole mouse brain. Thus, MouseBrainBench uses MICRONS as the main positive
empirical case for local structure-function analysis, while blocking
extrapolations to a complete mouse-brain digital twin.

\subsection{Simulation tools and mechanistic interpretation}

Several tools already exist for building and simulating biological neural
systems. BMTK and SONATA provide infrastructure for large-scale network models
\citep{gratiy2018bmtk,dai2020bmtk}, while The Virtual Brain provides a general
framework for connectome-based brain dynamics \citep{tvb2013}. These
systems are not replaced by MouseBrainBench. On the contrary, the proposal is
designed as an audit layer that can evaluate claims produced by models, data
adapters, and external simulation tools.

The central problem addressed in this paper is mechanistic identifiability. In
this context, a model is not considered mechanistically identifiable only because
it predicts well. It must also satisfy the evidence required by the claim:
reproducibility, anatomical specificity, directionality, structure-function
agreement, perturbation support when available, and controlled computational
cost. This distinction is the foundation of the proposed framework.
